% ============================================================================
% MUSTERLÖSUNG DS 3 - Gruppe A + B
% ============================================================================

\documentclass[11pt,a4paper]{article}

\usepackage[utf8]{inputenc}
\usepackage[T1]{fontenc}
\usepackage[ngerman]{babel}
\usepackage[left=2cm,right=2cm,top=1.8cm,bottom=1.5cm]{geometry}
\usepackage{helvet}
\usepackage{xcolor}
\usepackage{tabularx}
\usepackage{array}
\usepackage{enumitem}
\usepackage{tcolorbox}
\usepackage{fancyhdr}
\usepackage{lastpage}

\renewcommand{\familydefault}{\sfdefault}

\definecolor{spanisch}{HTML}{c41e3a}
\definecolor{grau}{HTML}{666666}
\definecolor{gruppeB}{HTML}{1565c0}
\definecolor{loesung}{HTML}{c62828}

\pagestyle{fancy}
\fancyhf{}
\fancyhead[L]{\small\textcolor{grau}{Spanisch Spätbeginner -- MUSTERLÖSUNG}}
\fancyhead[R]{\small\textcolor{grau}{ML-03}}
\fancyfoot[C]{\small\thepage/\pageref{LastPage}}
\renewcommand{\headrulewidth}{0.5pt}

\newcommand{\lsg}[1]{\textcolor{loesung}{\textbf{#1}}}

\newtcolorbox{gruppenbox}[1][]{
    colback=white,
    colframe=spanisch,
    fonttitle=\bfseries\color{white},
    colbacktitle=spanisch,
    title=#1,
    boxrule=1pt,
    arc=0pt,
    left=5pt, right=5pt, top=5pt, bottom=5pt
}

\newtcolorbox{gruppeBbox}[1][]{
    colback=white,
    colframe=gruppeB,
    fonttitle=\bfseries\color{white},
    colbacktitle=gruppeB,
    title=#1,
    boxrule=1pt,
    arc=0pt,
    left=5pt, right=5pt, top=5pt, bottom=5pt
}

\begin{document}

\begin{center}
{\Large\textbf{\textcolor{loesung}{MUSTERLÖSUNG}}}\\[0.3em]
{\large DS 3 -- Modalverben / Arbeiten im Ausland}
\end{center}

\vspace{0.5em}
\hrule
\vspace{1em}

% ============================================================================
% GRUPPE A
% ============================================================================
\begin{gruppenbox}[Gruppe A: Poder, querer, tener que (S. 32--33)]

\textbf{Aufgabe 1: Konjugationstabelle (9 P.)}

\begin{tabularx}{\textwidth}{@{}|c|l|l|l|@{}}
\hline
& \textbf{poder} & \textbf{querer} & \textbf{tener que} \\
\hline
yo & \lsg{puedo} & \lsg{quiero} & tengo que \\
\hline
tú & puedes & \lsg{quieres} & tienes que \\
\hline
él/ella & \lsg{puede} & quiere & \lsg{tiene que} \\
\hline
nosotros & podemos & \lsg{queremos} & tenemos que \\
\hline
vosotros & \lsg{podéis} & queréis & \lsg{tenéis que} \\
\hline
ellos & pueden & \lsg{quieren} & tienen que \\
\hline
\end{tabularx}

\vspace{0.5em}

\textbf{Aufgabe 2: Ergänzen (8 P.)}
\begin{enumerate}[label=\alph*), leftmargin=2em, nosep]
    \item ¿\lsg{Puedes} ir al cine esta noche?
    \item Nosotros \lsg{queremos} ir a la piscina.
    \item Ellos \lsg{tienen que} estudiar para el examen.
    \item Yo no \lsg{puedo} salir hoy.
    \item ¿Qué \lsg{queréis} hacer el sábado?
    \item María \lsg{tiene que} trabajar el domingo.
    \item ¿\lsg{Podemos} ir al parque después?
    \item Él \lsg{quiere} aprender español.
\end{enumerate}

\vspace{0.5em}

\textbf{Aufgabe 3: Zuordnung (4 P.)}

1 = \lsg{c}, 2 = \lsg{a}, 3 = \lsg{b}, 4 = \lsg{d}

\vspace{0.5em}

\textbf{Aufgabe 4: Sätze bilden (6 P.)}
\begin{enumerate}[label=\alph*), leftmargin=2em, nosep]
    \item \lsg{Yo quiero ir al cine.}
    \item \lsg{Nosotros podemos jugar al fútbol.}
    \item \lsg{Ella tiene que hacer los deberes.}
\end{enumerate}

\end{gruppenbox}

\vspace{1em}

% ============================================================================
% GRUPPE B
% ============================================================================
\begin{gruppeBbox}[Gruppe B: Trabajar en el extranjero (S. 105--106)]

\textbf{Aufgabe 1: Kasten ergänzen (4 P.)}

aprender idiomas = \lsg{Sprachen lernen}

conocer otras culturas = \lsg{andere Kulturen kennenlernen}

mejorar las oportunidades de trabajo = \lsg{Arbeitschancen verbessern}

desarrollar la independencia = \lsg{Selbstständigkeit entwickeln}

\vspace{0.5em}

\textbf{Aufgabe 2: Leseverstehen (8 P.) -- Beispielantworten}
\begin{enumerate}[label=\alph*), leftmargin=2em, nosep]
    \item \lsg{Carla trabaja en una empresa de marketing en Barcelona.}
    \item \lsg{Decidió ir a Barcelona para mejorar su español y ganar experiencia internacional.}
    \item \lsg{Jan hace prácticas en una empresa alemana en Madrid.}
    \item \lsg{Las ventajas son: aprende español, conoce la cultura española y mejora su currículum.}
\end{enumerate}

\vspace{0.5em}

\textbf{Aufgabe 3: Dilo con otras palabras (6 P.)}
\begin{enumerate}[label=\alph*), leftmargin=2em, nosep]
    \item \lsg{aprender cosas nuevas en el trabajo / trabajar y aprender al mismo tiempo}
    \item \lsg{saber más cosas / aprender más sobre un tema}
    \item \lsg{tener un mejor CV / tener más experiencias para encontrar trabajo}
\end{enumerate}

\vspace{0.5em}

\textbf{Aufgabe 4: Wörterbucharbeit (4 P.)}

las cualificaciones = \lsg{Qualifikationen}, el empleo = \lsg{Arbeitsstelle/Beschäftigung}

las prácticas = \lsg{Praktikum}, la estancia = \lsg{Aufenthalt}

\vspace{0.5em}

\textbf{Aufgabe 5: Schreibaufgabe (10 P.) -- Beispiel}

\lsg{Me gustaría trabajar en el extranjero para aprender idiomas y conocer otras culturas. Ya que hablo español, me gustaría ir a España o a Latinoamérica. Además, creo que es importante ganar experiencia internacional porque el mundo laboral es cada vez más global. También quiero desarrollar mi independencia y mejorar mis oportunidades de trabajo.}

\end{gruppeBbox}

\end{document}
